% !TeX root = ../../main.tex
\subsection{Propagation of the leading singularity}
\label{subsec:s08:02}

Fix a positive defect $h$ and $1/2<\nu<1$, using common domains for
all layers up to $h$ as in Corollary~\ref{crl:s08:domains}.
Put $S=R_\theta-x$. At $y=S$, the value of $\mathcal L_h$ below
means its continuous endpoint limit. The endpoint class of
Appendix~\ref{subsec:b:endpoints} defines unique actual coefficient
functions $P_h(\theta,x)$ by
\begin{equation}\label{eq:s08:leading-expansion}
 \begin{aligned}
 \mathcal L_h(\theta;x,y)
 ={}&\mathcal L_h(\theta;x,S)
 -\sqrt{\frac{2}{-w''(x_{*,\theta})}}
                    P_h(\theta,x)\sqrt{S-y}
 +E_h(\theta;x,y).
 \end{aligned}
\end{equation}
The remainder $E_h$ and the coefficient obey, uniformly on smaller
domains,
\begin{equation}\label{eq:s08:leading-bounds}
 \begin{aligned}
 |P_h(\theta,x)|&\leq C_h|S|^{(h-2)/2},\\
 |S^k\partial_y^k E_h(\theta;x,y)|
 &\leq C_{h,\nu,k}|S|^{(h-1)/2}|1-y/S|^{\nu-k}
                                      \quad(k=0,1,2).
 \end{aligned}
\end{equation}
All fixed higher or mixed derivatives are available after further
angular reduction, by Cauchy's formula. Uniqueness of the
square-root coefficient follows from $\nu>1/2$.
The negative-endpoint expansion is determined by
\eqref{eq:s07:borel-parity}.

\begin{lmm}[The logarithmic seed]\label{lmm:s08:seed}
On these inverse branches,
\begin{equation}\label{eq:s08:seed}
 P_1(\theta,x)=\frac1{x_{*,\theta}-v_\theta(x)}.
\end{equation}
\end{lmm}
\begin{prf}
In \eqref{eq:s07:borel-one}, only the positive $m=1$ term can
meet the critical value when $y\to S$. Its logarithm is
\[
 \log\frac{v_\theta(x+y)-v_\theta(x)}
                  {y\,\partial_xv_\theta(x)}.
\]
Insert \eqref{eq:s08:perturbed-inverse}. The numerator has
constant term $x_{*,\theta}-v_\theta(x)$ and square-root term
$-\sqrt{2/(-w''(x_{*,\theta}))}\sqrt{S-y}$.
The denominator is analytic and nonzero at the endpoint.
The square-root coefficient of the logarithm therefore gives
\eqref{eq:s08:seed}. The negative branch and $m\geq2$ terms are
analytic at this endpoint; their normal summability and the uniform
remainder follow from the first-layer argument in
Proposition~\ref{prp:b:common-domains}.
\end{prf}

The next calculation identifies the mechanism behind the higher
coefficients. The principal linear branch keeps its critical
distance fixed along its integration path. Characteristic
integration then turns its inverse-square-root derivative back
into a square root. The remaining branches have greater endpoint
regularity.

\begin{lmm}[Propagation of the principal coefficient]\label{lmm:s08:propagation}
For every fixed $h\geq2$,
\begin{equation}\label{eq:s08:theta-recurrence}
 P_h(\theta,x)
 =\frac1{x_{*,\theta}-v_\theta(x)}
                         \int_x^{R_\theta}P_{h-1}(\theta,s)\,\dd s.
\end{equation}
For complex centers the integral is continued along the straight
endpoint segment in the corresponding slit region.
\end{lmm}
\begin{prf}
\emph{The principal linear branch.}
In \eqref{eq:s07:linear-kernel}, its $m=1$ positive shift has
center $x'=x+(1-\tau)y$ and Borel argument $\tau y$.
Thus $R_\theta-x'-\tau y=S-y$ throughout the path.
Lemma~\ref{lmm:b:linear-endpoint} gives its square-root coefficient
by integrating the previous one at $y=S$; the coefficient is
\[
 -\sqrt{\frac{2}{-w''(x_{*,\theta})}}
            \int_0^1P_{h-1}(\theta,x+(1-\tau)S)\,\dd\tau.
\]
The endpoint bound for $P_{h-1}$ has exponent $(h-3)/2>-1$, so
the integral includes its critical end and is convergent there.

\smallskip\noindent
\emph{Characteristic recovery.}
The linear source is acted on by $1+y\partial_y$ before the
characteristic integral. Its differentiated square root has an
inverse-square-root singularity. The lower end of that integral
multiplies the original coefficient by
$S/(x_{*,\theta}-v_\theta(x))$, exactly as in
\eqref{eq:b:transport-endpoint-coefficient}.
Changing variable in the preceding $\tau$ integral gives
\eqref{eq:s08:theta-recurrence}.

To verify the analytic hypotheses of this calculation, the previous
layer has scale $(h-2)/2$ and remainder exponent $\nu$.
The other linear branches have a bounded first normalized derivative,
by Lemma~\ref{lmm:b:linear-endpoint}. The quadratic source has the
same scale and zero square-root coefficient, by
Lemma~\ref{lmm:b:quadratic-endpoint}: its two singular corners
cannot be reached together.

Both sources have order at least
$3h-2>(h-1)/2$ at $y=0$, as checked in
Proposition~\ref{prp:s07:borel-transport}. Thus
Lemma~\ref{lmm:b:local-transport} applies, including its uniform
remainder and derivative estimates.

The remote integral is analytic
at the endpoint by Corollary~\ref{crl:b:remote}; it contributes no
square-root coefficient. All operations are consequently valid on
the common complex regions, including their endpoint neighborhoods.
\end{prf}

For $\theta=0$ write $P_h(x)=P_h(0,x)$. In particular,
\begin{equation}\label{eq:s08:principal-recurrence}
 P_1(x)=\frac1{x_*-v_0(x)},\qquad
 P_h(x)=\frac1{x_*-v_0(x)}\int_x^R P_{h-1}(s)\,\dd s
                                      \quad(h\geq2).
\end{equation}
These coefficients are positive on the real interval $0\leq x<R$.
Their growth at the critical endpoint is controlled by
\eqref{eq:s08:leading-bounds}; the recurrence is therefore an
ordinary convergent integral identity, not a formal prescription.
